\subsection*{一步一步建立直观}

\paragraph{1. 先把上下文压缩成一次加权读取}

完整语言模型要根据此前的词元预测下一个词元。右图只抽出其中的注意力读取：两个候选信息各带一个分数，归一化后决定各自参与输出的份额。分数相同就平均读取。这些份额不是整个词表的下一词概率，二者在模型中的位置不同。



\paragraph{2. 相关性改变了信息的去向}

增加第二项相对分数，输出便更靠近它的值。可以把这一步理解为读者根据当前问题重新分配注意，而不是把所有上下文不加区分地平均。实际模型会从向量计算查询和键；本图直接调分数，用来隔离“分数如何影响权重”这一机制。



\paragraph{3. 集中读取不等于理解正确}

降低温度后，权重更加集中，但若较高分数本来分配给错误线索，集中只会强化错误。语言模型还需要词元表示、多层变换、训练目标及最终概率输出。这个窗口帮助理解一个局部部件，不把输出偏向某个值称为已经完成语言理解。
